\chapter{คู่มือการใช้งานแอปพลิเคชันและแถบเมนูวิชาการ}

\section{หน้าจอหลักและการตรวจจับสด (Live Camera Viewport)}
เมื่อเปิดแอปพลิเคชัน ระบบจะเริ่มต้นการทำงานของกล้องและโหลดโมเดล TensorFlow Lite ขึ้นสู่หน่วยความจำโดยอัตโนมัติ ที่มุมซ้ายบนจะแสดงสถานะการตรวจจับแบบสดและความเร็วเฟรมเรต เช่น \texttt{LIVE AI 30 FPS} ที่แถบควบคุมด้านบนขวา เกษตรกรสามารถแตะเพื่อเปิดโหมดกล่องมืด (Dark Chamber) เพื่อตัดแสงสะท้อนภายนอก ล็อกค่าแสง (Exposure Lock) หรือเปิดไฟฉายช่วยส่องใบไม้

\section{แถบเมนูควบคุมพื้นที่ตรวจจับ (Interactive ROI Floating Toolbar)}
ที่ด้านขวาของหน้าจอจะมีแถบลอยสำหรับควบคุมพื้นที่ตรวจจับ (ROI Floating Toolbar) ซึ่งได้รับการออกแบบตามมาตรฐานเดียวกับแอปพลิเคชันวิเคราะห์ดิน โดยประกอบด้วย 4 ส่วนสำคัญ

\subsection{การเลือกรูปทรงและปรับขนาดของกรอบตรวจจับ}
\begin{itemize}
  \item \textbf{กรอบสี่เหลี่ยม (Rectangle ROI)} เหมาะสำหรับการเล็งรอยโรคขนาดใหญ่หรือบริเวณแผ่นใบกว้าง
  \item \textbf{กรอบวงกลม (Circular ROI)} เหมาะสำหรับการตรวจวัดพื้นที่เนื้อเยื่อใบเขียวสมบูรณ์เพื่อประเมินค่า SPAD และ NPK
  \item \textbf{กรอบอิสระ (Freeform Polygon)} รองรับการวาดขอบเขตรอบรอยโรคที่มีรูปร่างไม่แน่นอน
  \item \textbf{ปุ่มขยายและย่อ (+ / -)} ปรับขนาดกรอบตรวจจับได้ตั้งแต่ $80$ ถึง $340$ พิกเซล
\end{itemize}

\subsection{แถบแสดงเฉดสีตามเกณฑ์มาตรฐานวิชาการ (Academic Color Strip)}
ด้านล่างของปุ่มควบคุมขนาด จะเป็นแถบเฉดสีแนวตั้งที่แสดงเกณฑ์มาตรฐานทางวิชาการ ผู้ใช้สามารถแตะปุ่มสลับโหมดเพื่อเปลี่ยนเกณฑ์ที่ต้องการติดตาม ได้แก่
\begin{enumerate}
  \item \textbf{โหมดคลอโรฟิลล์ (SPAD)} ไล่ระดับสี 5 ขั้น ตั้งแต่สีซีดขาว ($<25$) ไปจนถึงสีเขียวมรกตเข้ม ($>55$)
  \item \textbf{โหมดไนโตรเจน (N)} แสดงเกณฑ์ 5 ระดับ ตั้งแต่ขาดวิกฤต ($<1.8\%$) ไปจนถึงเกินเกณฑ์ ($>3.2\%$)
  \item \textbf{โหมดฟอสฟอรัส (P)} แสดงเกณฑ์ 4 ระดับ ตั้งแต่สีม่วงคล้ำ ($<0.12\%$) จนถึงระดับสูง
  \item \textbf{โหมดโพแทสเซียม (K)} แสดงแถบเตือนสีน้ำตาลขอบใบไหม้ ($<1.5\%$)
  \item \textbf{โหมดความรุนแรงของโรค} แสดงระดับการทำลายตั้งแต่ Grade 0 (ปกติ) จนถึง Grade 4 (ไหม้รุนแรง)
\end{enumerate}

\subsection{ไฟเรืองแสงนีออนและป้ายชี้เป้าสด (Neon Glow \& Live Pointer Badge)}
ขณะที่กล้องสตรีมภาพสดและ AI วิเคราะห์ค่าได้ แถบสีในช่องที่ตรงกับผลลัพธ์ปัจจุบันจะเปล่งแสงนีออนเรืองแสง พร้อมทั้งมีป้ายชี้เป้าสดทางด้านซ้าย เช่น \texttt{$\blacktriangleleft$ SPAD 48.5} หรือ \texttt{$\blacktriangleleft$ โรคราใบติด} เลื่อนไปยังตำแหน่งสีที่ถูกต้องแบบไดนามิก

\subsection{การเปิดแผ่นคำอธิบายเกณฑ์วิชาการ (Academic Legend Modal)}
เมื่อแตะที่แถบเฉดสีหรือไอคอนข้อมูล ระบบจะเปิดแผ่นหน้าต่าง Popover แสดงตารางเทียบเฉดสีตามเกณฑ์วิชาการ ค่าพิกเซลสีจริง ($RGB$ และ $CIE\ L^*a^*b^*$) พร้อมคำอธิบายทางสรีรวิทยาและการจัดการแปลง

\section{ขั้นตอนการสอนโมเดลเพิ่มเติม (Few-Shot Model Trainer)}
หากเกษตรกรพบอาการใบผิดปกติรูปแบบใหม่ สามารถแตะไอคอนรูปหมวกปริญญาด้านบน เพื่อเข้าสู่หน้าต่างสอนโมเดลเพิ่มเติม
\begin{enumerate}
  \item กรอกชื่อโรคหรืออาการเป็นภาษาไทยและภาษาอังกฤษ
  \item ระบุแนวทางการจัดการและรักษาในสวนทุเรียน
  \item กดปุ่มถ่ายภาพตัวอย่าง $2 - 3$ ใบ ระบบจะสกัดเวกเตอร์เอกลักษณ์ 128 มิติ
  \item กดปุ่ม \textbf{"บันทึกเข้าหน่วยความจำโมเดล"} ระบบจะจดจำและตรวจจับอาการดังกล่าวในกล้องสดได้ทันที
\end{enumerate}
